% Tait's reduction of vertex four-colouring to cubic edge three-colouring.
%
% Companion to foundations/map-euler-foundations.tex.  Every result carries
% an inline status marker; tools/check_status.py verifies the machine-checked
% ones against the Lean tree and generates the appendix ledger.

\section{The classical bridge: Tait's reduction}
\label{sec:tait-reduction}

The compositional argument produces a statement about \emph{cubic maps}: no
bridgeless spherical cubic map fails to be $3$-edge-colourable. The
Four-Colour Theorem is a statement about \emph{planar simple graphs}. Tait's
1880 reduction connects them, and this section gives it in full, with the
status of each ingredient marked, because the ingredients are not uniformly
available.

The reduction is classical and no part of it is in doubt. What is worth
setting down precisely is which pieces the present development already
supplies --- the algebraic heart is one of them --- and which remain to be
built, so that the remaining labour is visible rather than assumed.

\subsection{Statement}

\begin{definition}[Tait colouring]\label{def:tait}
A \emph{Tait colouring} of a cubic map is a proper $3$-edge-colouring: an
assignment of one of three colours to each edge such that the three edges at
every vertex receive three distinct colours.
\end{definition}

Throughout we identify the three colours with the three nonzero elements of
the Klein four-group
\[
  \mathbb{F}_2^2 \;=\; \{\,0,\ \mathsf a,\ \mathsf b,\ \mathsf a + \mathsf b\,\},
\]
which is legitimate because any bijection between colour sets will do, and
which buys the identity we shall use twice: \emph{the three nonzero elements
of $\mathbb{F}_2^2$ sum to zero}.

\begin{definition}[combinatorial planarity]
\label{def:combinatorial-planarity}
\Sverified{Mettapedia.GraphTheory.FourColor.GoertzelV24SphericalGraphPresentation.CombinatoriallyPlanar}
A finite simple graph is \emph{combinatorially planar} when each connected
component with at least three vertices carries a connected cellular spherical
rotation-system presentation.  Thus the darts over every vertex have one
cyclic order, the primal graph is the given component, and the face orbits
satisfy Euler's spherical equation.
\end{definition}

\begin{theorem}[Tait's reduction for finite combinatorial maps]
\label{thm:tait-reduction}
\Sverified{Mettapedia.GraphTheory.FourColor.GoertzelV24StellarTaitReduction.combinatorialFourColorStatement_of_tait}
Suppose every bridgeless spherical cubic map admits a Tait colouring. Then
every finite combinatorially planar simple graph admits a proper vertex
colouring with four colours.
\end{theorem}

\begin{corollary}[topological-drawing formulation]
\label{cor:tait-reduction-topological}
\Sconditional{Lemma~\ref{lem:planar-spherical-presentation}}
Suppose every bridgeless spherical cubic map admits a Tait colouring. Then
every finite planar simple graph, with planarity defined by a crossing-free
topological drawing in the plane, admits a proper vertex colouring with four
colours.
\end{corollary}

\begin{proof}
Apply Lemma~\ref{lem:planar-spherical-presentation} componentwise and then
Theorem~\ref{thm:tait-reduction}.
\end{proof}

\begin{proposition}[audit of the legacy Lean planarity predicate]
\label{prop:weak-planarity-refuted}
\Sverified{Mettapedia.GraphTheory.FourColor.GoertzelV24WeakPlanarityRefutation.current_planarity_does_not_imply_four_colorable}
The predicate currently named \texttt{Mettapedia.GraphTheory.IsPlanar} is not
a formalization of planarity.  Every finite simple graph satisfies it, while
$K_5$ satisfies it and is not $4$-colourable.
\end{proposition}

\begin{proof}
The current structure records only a finite family of edge sets in which
every edge occurs twice.  For an arbitrary graph take two formal faces and
declare the boundary of each to be the entire edge set.  This satisfies every
field.  On the other hand, a colouring of the complete graph on five vertices
is injective, so no map from its vertices to a four-element palette can be a
proper colouring.  Both statements, and their conjunction, are checked by
the declaration named above.
\end{proof}

This proposition concerns only the repository's legacy vocabulary, not
Theorem~\ref{thm:tait-reduction}.  The controlling formal statement uses
the genuine combinatorial notion already used throughout the route: a
graph-backed rotation system with cyclic vertex rotations, connected primal
graph, and spherical Euler equation.  The stellar construction below then
acts on that presentation.  No theorem may use the legacy
\texttt{IsPlanar} premise as a substitute, since that would make the formal
headline false before the compositional argument begins.

\begin{lemma}[the abstract-planar presentation bridge]
\label{lem:planar-spherical-presentation}
\Sconditional{the classical finite plane-embedding representation theorem}
Every finite connected planar simple graph with at least three vertices
admits a connected spherical graph presentation: a dart rotation system
whose primal graph is the given graph, whose rotation is cyclic at every
vertex, and whose face orbits satisfy
\[
  |V|-|E|+|F|=2.
\]
\end{lemma}

\begin{proof}
Choose a crossing-free embedding of the abstract planar graph in the sphere.
Take the two oriented incidences of every edge as darts, reverse a dart to
obtain the fixed-point-free edge involution, and read the cyclic order of
the incident darts around each vertex from the orientation of the sphere.
The resulting primal graph is the original graph.  Its face permutation
traces the complementary regions of the embedding, and Euler's formula gives
the displayed identity.  This is the standard equivalence between a cellular
sphere embedding and a spherical rotation system.

The marker is \textsc{conditional}, rather than \textsc{prose}, because the
repository currently has no machine-checked constructor implementing this
topological representation theorem.  The ordinary argument above fixes the
intended bridge exactly; it must not be replaced by the refuted legacy
predicate of Proposition~\ref{prop:weak-planarity-refuted}.  This lemma is the
sole extra input for Corollary~\ref{cor:tait-reduction-topological}; it is not
a premise of the combinatorial-map Four-Colour statement in
Theorem~\ref{thm:tait-reduction}.
\end{proof}

The proof occupies the rest of the section. It is assembled from five
ingredients, stated separately so that each carries its own status.

\subsection{The ingredients}

\begin{lemma}[componentwise spherical reduction]\label{lem:components}
\Sverified{Mettapedia.GraphTheory.FourColor.GoertzelV24SphericalGraphPresentation.sphericalFourColorStatement_of_connected}
Suppose every connected graph on at least three vertices which carries a
connected spherical graph presentation is four-colourable.  Then every
finite graph whose nontrivial connected components carry such presentations
is four-colourable.
\end{lemma}

\begin{proof}
A component on at most two vertices is coloured injectively into the
four-element palette.  Apply the connected hypothesis to every other
component.  The connected components partition the vertex set, and every
edge has both endpoints in one component, so the component colourings combine
to one proper colouring of the whole graph.  The abstract-planar input is
kept separate: Lemma~\ref{lem:planar-spherical-presentation} supplies the
required presentation to each connected planar component.
\end{proof}

\begin{lemma}[triangulation]\label{lem:triangulate}
\Sopen
Every connected planar simple graph $G$ on at least three vertices is a
spanning subgraph of a plane simple graph $T$ on the same vertex set in
which every face is bounded by a triangle.
\end{lemma}

This is the standard maximality argument: fix a plane embedding and add
non-crossing edges between non-adjacent vertices as long as any can be
added; the process terminates because the graph is simple and finite, and a
maximal plane simple graph on at least three vertices is a triangulation. It
is classical, and it is \emph{not} formalized here. Of the ingredients it is
the one carrying real formalization cost, since it needs a plane embedding
to be manipulated rather than merely assumed. It is also avoidable:
Theorem~\ref{thm:stellar-darts} builds a triangulation outright, on darts,
and the reduction can be routed through that instead --- at the price of
letting the vertex set grow, which the reduction does not mind. This lemma
is retained in its classical form, and keeps its \textsc{open} marker,
because that is the statement it makes; the reroute does not prove it.

Because a proper colouring of $T$ restricts to a proper colouring of any
spanning subgraph, it suffices to colour $T$.

\subsection{Duality, at the level the reduction uses it}

The dual is usually introduced by drawing a vertex inside each face and an
edge across each primal edge, and that picture makes duality look like a
statement about embeddings. At the level this reduction needs, it is not.
It is a one-line change of generators.

\begin{definition}[the dual map]\label{def:dual-map}
\Sprose
Let $M = (D, \sigma, \alpha)$ be a map, with $\varphi = \sigma\alpha$ its
face permutation. Its \emph{dual} is
\[
  M^{*} \;=\; (D,\; \varphi,\; \alpha).
\]
The darts and the edge flip are unchanged; only the vertex rotation is
replaced by the face permutation.
\end{definition}

\begin{lemma}[what the dual inherits, for free]\label{prop:dual-map}
\Sverified{Mettapedia.GraphTheory.FourColor.GoertzelV24DualMap.dualRotation_mul_flip}
\Sverified{Mettapedia.GraphTheory.FourColor.GoertzelV24DualMap.orbitCount_dualRotation_mul_flip}
\Sverified{Mettapedia.GraphTheory.FourColor.GoertzelV24DualMap.wordOrbitCount_dualRotation}
\Sverified{Mettapedia.GraphTheory.FourColor.GoertzelV24DualMap.orbitCount_dualRotation_add}
\Sverified{Mettapedia.GraphTheory.FourColor.GoertzelV24DualMap.dualRotation_dualRotation}
Let $M = (D,\sigma,\alpha)$ be a map. Then:
\begin{enumerate}
\item $M^{*}$ is a map: $\alpha$ is unchanged, hence still a fixed-point-free
  involution, and $\varphi$ is a permutation of $D$.
\item The face permutation of $M^{*}$ is $\varphi\alpha = \sigma\alpha\alpha
  = \sigma$. So the faces of the dual are the vertices of the primal:
  $V^{*} = F$, $E^{*} = E$, $F^{*} = V$.
\item $\langle \varphi, \alpha\rangle = \langle \sigma, \alpha\rangle$, since
  $\varphi = \sigma\alpha$ and $\sigma = \varphi\alpha$. Hence $M$ and $M^{*}$
  have the same orbits on darts and the same number of components; one is
  connected exactly when the other is.
\item $V^{*} - E^{*} + F^{*} = F - E + V$, so $M^{*}$ is spherical exactly
  when $M$ is.
\item $M^{**} = M$ on the nose: the dual of $(D,\varphi,\alpha)$ is
  $(D, \varphi\alpha, \alpha) = (D,\sigma,\alpha)$.
\end{enumerate}
\end{lemma}

\begin{proof}
Every item is the displayed computation. Item (v) is worth remarking on:
duality here is an involution up to \emph{equality}, not up to isomorphism,
which is what makes it usable without transport lemmas.
\end{proof}

\begin{corollary}[loops and bridges exchange]\label{cor:loop-bridge}
\Sverified{Mettapedia.GraphTheory.FourColor.GoertzelV24DualMap.two_le_wordOrbitCount_dual_of_loop}
\Sverified{Mettapedia.GraphTheory.FourColor.GoertzelV24DualMap.sameCycle_of_not_wordReachable_dual}
\Sverified{Mettapedia.GraphTheory.FourColor.GoertzelV24DualMap.wordReachable_dual_of_not_loop}
Let $M$ be a connected spherical map.  An edge $e$ is a loop of $M$ if and
only if it is a bridge of $M^{*}$.  Consequently the dual of a connected
spherical loopless map is bridgeless.
\end{corollary}

\begin{proof}
Write $d, \alpha d$ for the darts of $e$. That $e$ is a loop of $M$ means its
two darts sit at one vertex, that is, in one $\sigma$-orbit. By
Lemma~\ref{prop:dual-map}(ii) the $\sigma$-orbits are exactly the face
orbits of $M^{*}$. So $e$ is a loop of $M$ iff its two darts lie on a common
face of $M^{*}$, and by the two directions of the bridge characterisation
--- Theorems~\ref{thm:same-face-bridge} and~\ref{thm:converse-packaged},
both applied to the map $M^{*}$ --- that holds iff $e$ is a bridge of
$M^{*}$.
\end{proof}

\begin{remark}[what this changes in the accounting]
\Sprose
Nothing in the last two results mentions a plane, a curve, or a face
interior. Duality, the exchange of loops with bridges, and the preservation
of sphericity and connectedness are all consequences of replacing $\sigma$ by
$\sigma\alpha$. So of the embedding obligations this section owes, the
\emph{duality} one is not an embedding obligation: it is elementary
permutation algebra sitting on top of a bridge characterisation that is
already machine-checked.  Stellar subdivision below supplies the only
triangulation the reduction consumes, and the graph-backed argument of
Section~\ref{sec:graph-backed} has a combinatorial digon proof.  Thus no
separation principle remains in the controlling ordinary-mathematics route.
\end{remark}

\begin{lemma}[the dual of a triangulation is cubic and bridgeless]\label{lem:dual-cubic}
\Sprose
Let $T$ be a connected spherical loopless map whose face orbits have length
three and meet three distinct vertices. Then its dual map $T^{*}$ is connected,
spherical, cubic, and bridgeless.
\end{lemma}

\begin{proof}
By Lemma~\ref{prop:dual-map}, connectedness and sphericity pass to the dual.
Each face orbit of $T$ has three darts, so each vertex rotation of $T^{*}$ has
three darts: $T^{*}$ is cubic.  The distinct-vertex hypothesis records that
these are genuine triangular faces; it will be supplied by stellar
subdivision.

Bridgelessness is subtler than it first appears, and it is worth stating
exactly which implication is needed. Since $T$ is loopless,
what must be shown is
\[
  \text{no primal loop} \;\Longrightarrow\; \text{no dual bridge},
  \qquad\text{equivalently}\qquad
  \text{dual bridge} \;\Longrightarrow\; \text{primal loop}.
\]
Now observe that the dual of the map $(\sigma, \alpha)$ is the map
$(\varphi, \alpha)$ with $\varphi = \sigma\alpha$: its faces are the orbits
of $\varphi\alpha = \sigma$, that is, the primal vertices. Consequently a
\emph{primal loop} is exactly an edge of the dual whose two sides lie on one
dual face, and $V' + F' = F + V = E + 2$, so the dual of a spherical map is
spherical with the same edge count.

Corollary~\ref{cor:loop-bridge} now supplies exactly this: an edge is a loop
of $T$ precisely when it is a bridge of $T^{*}$, so the loopless primal has a
bridgeless dual, with no appeal to an embedding.
\end{proof}

\begin{corollary}[the explicit stellar dual has the complete Tait class]
\label{cor:stellar-dual-package}
\Sverified{Mettapedia.GraphTheory.FourColor.GoertzelV24StellarDualStructure.stellarDual_bridgelessSphericalCubicMapData}
Let $M$ be a connected spherical graph presentation with at least three old
vertices, let $T$ be its stellar subdivision, and let $T^{*}$ be the facial
dual.  Then $T^{*}$ is a connected, bridgeless, spherical cubic rotation map
with one cyclic rotation at every vertex.
\end{corollary}

\begin{proof}
The stellar map is connected and spherical, every stellar face orbit has
three darts, and every stellar edge has two distinct face sides.  Replacing
the stellar vertex rotation by its face permutation therefore makes the dual
cubic and preserves connectedness and Euler characteristic.  If a dual edge
were a bridge, the two darts of the corresponding stellar edge would lie in
one dual face orbit; dual faces are stellar vertex orbits, so that stellar
edge would be a loop, contrary to the looplessness of the stellar
construction.  The declaration above packages these facts as exactly the
map class consumed by the Tait-colouring hypothesis.
\end{proof}

\begin{remark}[a correction worth recording]\label{rem:not-one-interface}
\Sprose
It is tempting, having listed the residual obligations of this section and
of Section~\ref{sec:graph-backed}, to conclude that they are all facets of a
single plane-duality interface. That is very nearly right and is wrong in
one place. Adding a non-crossing edge, maximality forcing triangles, and a
digon's parallel pair bounding a face are indeed three statements about
plane embeddings and their duals. The bridgelessness of $T^{*}$ is not: by
the duality just described it is a statement about permutations, namely the
converse of Theorem~\ref{thm:same-face-bridge}, and it requires no embedding
machinery at all.

The distinction is worth keeping because the two kinds of debt have very
different costs. The converse is an isolated combinatorial lemma --- now proved, by a walk
argument that never counts anything; the embedding interface is a
development.
Merging them in the ledger would have hidden a cheap item inside an
expensive one.
\end{remark}

We record the fact that makes the algebra work.

\begin{lemma}[the Klein identity]\label{lem:klein}
\Sverified{Mettapedia.GraphTheory.FourColor.GoertzelV24TaitReductionPotential.klein_sum_eq_zero}
\Sverified{Mettapedia.GraphTheory.FourColor.GoertzelV24TaitReductionPotential.klein_exhausts}
The three nonzero elements of $\mathbb{F}_2^2$ sum to zero, and any three
pairwise distinct nonzero elements of $\mathbb{F}_2^2$ are all three of them.
\end{lemma}

\begin{proof}
$\mathsf a + \mathsf b + (\mathsf a + \mathsf b) = 0$ since the group has
exponent two. There are exactly three nonzero elements, so three pairwise
distinct ones exhaust them.
\end{proof}

\begin{theorem}[facial boundaries span the cycle space]\label{thm:cycle-space}
\Sverified{Mettapedia.GraphTheory.FourColor.GoertzelV24OrbitFaceCycleSpaceEquality.range_orbitFaceBoundaryLinearMap_eq_f2CycleSpace}
\Sverified{Mettapedia.GraphTheory.FourColor.GoertzelV24RotationCycleSpace.range_faceBoundaryLinearMap_eq_cycleSpace}
For a connected spherical map with two-sided faces and connected dual, the
image of the facial-boundary map is exactly the $\mathbb{F}_2$ cycle space of
the literal edge carrier of the map.
\end{theorem}

This is the algebraic heart of the reduction and it is already available.
Note its hypothesis of two-sided faces: that is supplied by
Corollary~\ref{cor:two-sided} of the previous section, so the two halves of
this development meet here.  The second checked declaration is the form
consumed below.  It uses the map's edge type as the columns of the incidence
matrix, so two parallel edges remain two different coordinates.  This is
load-bearing for stellar subdivision: if an old face revisits an old vertex,
two different spokes may have the same pair of endpoints, and projecting to
a simple-graph edge set would identify them.

\subsection{Proof of the reduction}

\begin{proof}[Proof of Theorem~\ref{thm:tait-reduction}]
By Lemma~\ref{lem:components} we may assume $G$ is connected with at least
three vertices.  Let $M$ be the spherical rotation-system presentation
supplied by Definition~\ref{def:combinatorial-planarity}; it is connected and
loopless because its primal graph is the connected simple graph $G$.  Apply
Definition~\ref{def:stellar} to obtain $T=M'$.  By
Theorem~\ref{thm:stellar-darts}, $T$ is connected, spherical and loopless, it
contains every old vertex and edge of $G$, and every face is a $3$-cycle
meeting three distinct vertices; the same theorem also gives two distinct
face sides at every edge.  Hence a proper colouring of $T$ restricts to one
of $G$, and the two-sided hypothesis used in Step~2 below is available.

By Lemma~\ref{lem:dual-cubic} the dual $T^{*}$ is a bridgeless spherical
cubic map, so the hypothesis of the theorem furnishes a Tait colouring of
$T^{*}$. Primal and dual share an edge set, so this is a function
\[
  \omega \colon E(T) \longrightarrow \mathbb{F}_2^2 \setminus \{0\}
\]
with the property that the three edges bounding any face of $T$ --- equally,
the three edges at the corresponding vertex of $T^{*}$ --- receive three
pairwise distinct nonzero values.

\emph{Step 1: $\omega$ vanishes on facial boundaries.} Fix a face of $T$. Its
three edges receive three pairwise distinct nonzero values, which by
Lemma~\ref{lem:klein} are all three nonzero elements, and these sum to zero.
So $\omega$ sums to zero over the boundary of every face.

\emph{Step 2: $\omega$ vanishes on the whole cycle space.} By
Theorem~\ref{thm:cycle-space} the facial boundaries span the $\mathbb{F}_2$
cycle space of $T$. Summation of $\omega$ over an edge set is
$\mathbb{F}_2$-linear in that set, applied coordinatewise in
$\mathbb{F}_2^2$; a linear functional vanishing on a spanning set vanishes
identically. So $\omega$ sums to zero over every element of the cycle space.

\emph{Step 3: potentials exist.} Fix a root $r \in V(T)$. For $v \in V(T)$
choose a path $P_{r,v}$ from $r$ to $v$ --- one exists by connectedness ---
and set
\[
  c(v) \;:=\; \sum_{e \in P_{r,v}} \omega(e) \;\in\; \mathbb{F}_2^2 .
\]
This is well defined: if $P$ and $P'$ are two $r$--$v$ paths, their symmetric
difference is an element of the cycle space, so by Step~2 the two sums agree.
More generally, the same value is obtained from any $r$--$v$ walk: retain the
edges traversed an odd number of times (even repetitions cancel in
$\mathbb F_2^2$); its symmetric difference with $P_{r,v}$ has even degree at
every vertex and therefore lies in the cycle space.

\emph{Step 4: the potential is a proper colouring.} Let $uv \in E(T)$.
Appending the edge $uv$ to a path from $r$ to $u$ gives a walk from $r$ to
$v$, so the preceding walk-independence gives
$c(v) = c(u) + \omega(uv)$, that is
\[
  c(u) + c(v) \;=\; \omega(uv) \;\neq\; 0 ,
\]
so $c(u) \neq c(v)$. Thus $c \colon V(T) \to \mathbb{F}_2^2$ is a proper
colouring with four colours, and its restriction to $G$ is one too.
\end{proof}

The algebraic implication used in Steps~1--4---from a facial Tait colouring,
two-sidedness, connected primal and dual, and spherical cubic map data to a
proper four-colouring---is now machine-checked both on the earlier
graph-backed carrier and on the literal multigraph carrier used by stellar
subdivision.
\Sverified{Mettapedia.GraphTheory.FourColor.GoertzelV24TaitReductionPotential.colorable_four_of_facialTaitColouring}
\Sverified{Mettapedia.GraphTheory.FourColor.GoertzelV24RotationWalkCycleParity.incidenceMatrix_mulVec_walkEdgeParity}
\Sverified{Mettapedia.GraphTheory.FourColor.GoertzelV24RotationTaitPotential.colorable_four_of_dualTaitColoring}

The final consumer is checked as well.  A Tait colouring of the explicit
stellar facial dual is integrated on literal stellar edges; the resulting
four-colouring of the stellar primal map is then restricted along the old
vertex inclusion.  Thus the entire reverse Tait implication is verified for
the genuine connected spherical presentation used in this section.
\Sverified{Mettapedia.GraphTheory.FourColor.GoertzelV24StellarTaitReduction.stellar_adj_of_adj}
\Sverified{Mettapedia.GraphTheory.FourColor.GoertzelV24StellarTaitReduction.colorable_four_of_stellarDual_taitColorable}
\Sverified{Mettapedia.GraphTheory.FourColor.GoertzelV24StellarTaitReduction.connectedSphericalFourColorStatement_of_tait}
\Sverified{Mettapedia.GraphTheory.FourColor.GoertzelV24StellarTaitReduction.sphericalFourColorStatement_of_tait}

\begin{proposition}[exact entry into zero Count]
\label{prop:stellar-zero-count-entry}
\Sverified{Mettapedia.GraphTheory.FourColor.GoertzelV24RotationTaitCount.rotationSystemTaitCount_eq_zero_iff}
\Sverified{Mettapedia.GraphTheory.FourColor.GoertzelV24StellarTaitReduction.stellarDual_taitCount_eq_zero_of_not_colorable}
For a finite closed rotation system \(R\), let

\[
 \operatorname{Count}_{\mathrm{Tait}}(R)
 :=\#\{C:C\text{ is a proper nonzero Tait edge colouring of }R\}.
\]
Then
\[
 \operatorname{Count}_{\mathrm{Tait}}(R)=0
 \quad\Longleftrightarrow\quad
 R\text{ is not Tait-colourable}.
\]
Moreover, if a connected spherical presentation \(M\) with at least three
vertices is not four-colourable, then its explicit stellar facial dual
\(T^{*}\) has zero Tait Count.
\end{proposition}

\begin{proof}
The first assertion is finite cardinality: the displayed count is positive
exactly when its finite type of witnesses is inhabited, and inhabitation is
the definition of Tait-colourability.  For the second assertion, suppose
instead that \(T^{*}\) had positive Count.  Choose the corresponding Tait
colouring.  The checked Klein-group integration construction gives a proper
four-colouring of the stellar primal map, and restriction to the old vertices
gives a proper four-colouring of \(M\), a contradiction.

This is the exact obstruction transport required by the reductive spine.  It
does not claim the converse local statement that every four-colouring of the
old graph extends over all stellar face centres; that stronger statement is
neither needed nor generally supplied by stellar subdivision.
\end{proof}

\begin{theorem}[the sound M0 least-counterexample wrapper]
\label{thm:spherical-headline-from-no-minimal}
\Sverified{Mettapedia.GraphTheory.FourColor.GoertzelV24SphericalMinimalCounterexampleSelection.everyBridgelessSphericalCubicTaitColorable_of_no_minimal}
\Sverified{Mettapedia.GraphTheory.FourColor.GoertzelV24SphericalMinimalCounterexampleSelection.sphericalFourColorStatement_of_no_minimal}
\Sverified{Mettapedia.GraphTheory.FourColor.GoertzelV24SphericalMinimalCounterexampleSelection.not_noGraphBackedVertexMinimalTaitCounterexample_of_not_colorable}
Suppose there is no graph-backed vertex-minimal non-Tait-colourable member of
the connected, bridgeless, spherical cubic rotation-map class.  Then every
bridgeless spherical cubic rotation map is Tait-colourable, and every finite
graph carrying spherical presentations on its nontrivial connected
components is four-colourable.
\end{theorem}

\begin{proof}
If a non-Tait-colourable map existed, choose one with the least number of
vertices.  Corollary~\ref{cor:two-sided} constructs its two-sided face
property from bridge-freeness; the digon normalization of
Section~\ref{sec:graph-backed} excludes parallel edges by minimality.  The
canonical primal simple graph therefore supplies the graph-backed
vertex-minimal counterexample forbidden by the hypothesis.  Hence every map
in the class is Tait-colourable.  The checked stellar construction and
Klein-group integration above give the connected spherical Four-Colour
statement, and Lemma~\ref{lem:components} combines the component colourings.

This well-ordering wrapper is important logically: two-sidedness and graph
backing are conclusions inside the least-counterexample argument, not extra
premises exported to the compositional core.
\end{proof}

\subsection{What the reduction costs}

It is worth being explicit about where the remaining labour sits.  The
reverse Tait implication is now complete both as ordinary mathematics and as
a formal theorem over connected spherical graph presentations.

Steps~1 through~4 are machine-checked algebra over $\mathbb{F}_2^2$ resting
on Theorem~\ref{thm:cycle-space}.  The other ingredients are now
equally explicit: Definition~\ref{def:stellar} and
Theorem~\ref{thm:stellar-darts} replace maximality by a finite permutation
construction; Definition~\ref{def:dual-map} is a change of generator; and
Corollary~\ref{cor:loop-bridge} supplies dual bridgelessness.  The
componentwise wrapper is checked by Lemma~\ref{lem:components}, so the finite
combinatorial-map form of Tait's reduction has no remaining premise besides
the Tait-colourability hypothesis it is supposed to reduce from.  The named
representation bridge of Lemma~\ref{lem:planar-spherical-presentation} is
needed only to translate the separate topological-drawing convention of
Corollary~\ref{cor:tait-reduction-topological}.  No legacy
\texttt{IsPlanar} premise is used.

\subsection{Audit of the unused maximality route}

The next discussion records why the textbook maximal-plane-graph route was
not chosen.  It is retained as an audit of an attractive detour, not as a
dependency of Theorem~\ref{thm:tait-reduction}; the controlling proof uses
stellar subdivision instead.

\paragraph{Adding a non-crossing edge is a construction, not a theorem.}
Let $d, d'$ be darts on one face orbit. Adjoin two darts $x, y$; set
$\alpha' = \alpha \cup (x\,y)$; and splice $x$ into the rotation at $d$'s
vertex and $y$ into the rotation at $d'$'s vertex, so $\sigma' d = x$,
$\sigma' x = \sigma d$, and symmetrically at $d'$. No embedding is consulted:
the rotation system \emph{is} the embedding. The Euler bookkeeping is
immediate, since $V' = V$, $E' = E + 1$, and the face through $d, d'$ splits,
so $F' = F + 1$ and $V' - E' + F' = V - E + F$. Sphericity is preserved, and
the face count is governed by the split law of
Lemma~\ref{lem:split-merge} rather than by a separation argument.

\paragraph{Maximality is where the content actually sits.}
The step that resists is not the construction but the claim that a maximal
simple plane graph on at least three vertices has only triangular faces.
Given a face of length at least four, one wants two of its vertices that are
\emph{not already adjacent elsewhere in the graph}, since joining adjacent
ones would destroy simplicity. That non-adjacent pair must be shown to exist,
and its existence is a genuine combinatorial fact about plane simple graphs,
not a corollary of the construction above.

\paragraph{Why counting does not finish it.}
It is tempting to hope Euler settles maximality directly, as it settled
two-sidedness. It does not, and the gap is instructive. If a simple plane
graph with $V \ge 3$ has one face of length at least four and the rest at
least three, then $2E \ge 3F + 1$, and substituting $F = 2 - V + E$ gives
$E \le 3V - 7$. So the graph falls short of the triangulation bound
$E = 3V - 6$. But that shows only that it is not edge-\emph{maximum}; being
\emph{maximal}, in the sense that no further edge can be added, is a weaker
condition. Counting establishes a deficit without exhibiting an edge to add,
and producing that edge is what needs an argument about which vertices the
face separates.

\begin{remark}[where the last classical debt lands]
\Sprose
Within the unused maximality route, combinatorial separation for rotation
systems is available in dual form below.  An earlier draft then claimed that,
on a face of length at least four, a chord through the first and third
vertices separates the second vertex from the fourth.  That inference is
wrong: the second vertex lies on the resulting three-cycle.  Recovering the
maximality route would require the additional rotation-at-the-second-vertex
argument discussed below; the controlling stellar route needs none of it.

This is retained only as an audit of a discarded proof route.  The generic
dual separation result is useful elsewhere, but it does not by itself prove
the vertex non-adjacency asserted by the former maximality argument.
\end{remark}

\paragraph{Separation may already be available.}
Before treating combinatorial separation as a development to be built, it is
worth noticing that the pieces for it are in hand. Let $C$ be a cycle. It
lies in the $\mathbb{F}_2$ cycle space, so by Theorem~\ref{thm:cycle-space}
it is a sum of facial boundaries, say $C = \sum_{f \in S} \partial f$ for
some set $S$ of faces. By Corollary~\ref{cor:two-sided} each edge has two
\emph{distinct} faces, so $e \in \partial f$ exactly when $f$ is one of them,
and therefore
\[
  e \in C \quad\Longleftrightarrow\quad
  \bigl|\{\text{the two faces of } e\} \cap S\bigr| \text{ is odd,
  that is, exactly one.}
\]
Writing $\mathrm{side}(f) := (f \in S)$, an edge outside $C$ has its two
faces on the same side and an edge of $C$ has them on opposite sides. That
is separation: $C$ is precisely the edge cut between the two classes, and
nothing outside $C$ crosses.

\begin{theorem}[the boundary map is the two-sided sum]\label{thm:pair-sum}
\Sverified{Mettapedia.GraphTheory.FourColor.GoertzelV24FaceSeparation.orbitFaceBoundaryLinearMap_apply_eq_pair_sum}
With two-sided faces, the boundary map at an edge returns the sum of the
coefficients of that edge's two faces.
\end{theorem}

\begin{corollary}[separation, dual form]\label{cor:separation}
\Sverified{Mettapedia.GraphTheory.FourColor.GoertzelV24FaceSeparation.mem_image_iff_separates}
An edge lies in the image of a face-set indicator exactly when its two sides
fall on opposite sides of the set. Hence the image is the edge cut between
that set and its complement, and no edge outside the cut crosses it.
\end{corollary}

With Theorem~\ref{thm:cycle-space}, every cycle is such a cut: the
combinatorial separation that planar arguments normally import from topology
is available here without importing anything.

\begin{remark}[what remains of the maximality step, and two routes to it]
\Sprose
Corollary~\ref{cor:separation} separates \emph{faces}. The maximality
argument consumes a statement about \emph{vertices}, so a transfer is still
needed. Since that transfer is now the only unwritten step between this
section and a complete classical bridge, it is worth setting down the two
candidate routes precisely, neither of which is established here or consumed
by the controlling proof.

Suppose, for contradiction, that every two vertices on a face of length
$k \ge 4$ are adjacent.

\emph{Route A, through separation.} The cycle $C$ formed by two face edges
$v_1v_2$, $v_2v_3$ and the chord $v_3v_1$ is a cycle, so by
Corollary~\ref{cor:separation} it is a face cut. But $v_2$ lies \emph{on}
$C$, while the conclusion wanted concerns $v_2$ and $v_4$ being
non-adjacent rather than lying on opposite sides. The step from a face
two-colouring to a vertex non-adjacency is not immediate, and no transfer
supplying it has been identified.

\emph{Route B, through forbidden subgraphs.} If $k \ge 5$ the face's
vertices induce $K_5$, which is not planar. If $k = 4$ they induce $K_4$
with all four vertices on one face; but every face of $G$ lies inside a face
of the subgraph $K_4$, and each of those is bounded by only three vertices,
so no face of $G$ can meet all four. This route needs the non-planarity of
$K_5$ and the containment of faces under passage to a subgraph --- both
classical, neither present in the current libraries.

Route~B does not use Corollary~\ref{cor:separation} at all, which is worth
noticing on its own: the separation result may be valuable elsewhere in this
development without being the tool this particular step wants.

It is tempting to call Route~B the shorter, and that judgement does not
survive checking. No current library supplies planarity in any form, so the
non-planarity of $K_5$ would have to be proved here, and it does not follow
from the Euler bound of Theorem~\ref{thm:euler-bound}: the classical
derivation $2 = V - E + F$ together with $F \le 2E/3$ yields $E \le 3V-6$,
but what is available here is the \emph{inequality} $V - E + F \le 2C$, and
two upper bounds in the same direction do not combine to bound $E$. Route~B
therefore needs sphericity as a theorem about embeddings rather than as the
hypothesis it presently is.

Both routes accordingly require real work, and no preference between them is
recorded here. The estimate that one was shorter is withdrawn.
\end{remark}

\begin{remark}[a third route, which avoids maximality entirely]\label{rem:stellar}
\Sprose
Both routes above exist only to serve Lemma~\ref{lem:triangulate}, and that
lemma was stated in the classical way --- extend $G$ to a \emph{maximal}
plane graph --- because that is how the textbooks do it. Maximality is not
what the reduction needs. What it needs is \emph{some} plane triangulation
containing $G$, and one can be built outright.

Place a new vertex $u_F$ inside each face $F$ and join it to every vertex on
$F$'s boundary. A face of length $k$ becomes $k$ triangles. Nothing is
removed, so $G$ survives as a subgraph; a proper colouring of the result
restricts to $G$, since restricting a proper colouring to a subset of
vertices is again proper; every face is a triangle, so the dual is cubic;
and planarity is immediate because $u_F$ sits inside $F$.

Crucially, \emph{no edge is added between existing vertices}. The entire
difficulty of maximality was finding two non-adjacent vertices on a long
face, and this construction never needs one.

It is tempting to attach a caveat here --- that the construction yields a
\emph{simple} graph only when each face boundary has distinct vertices,
hence only for $2$-connected $G$ --- and to trade maximality for a reduction
to blocks. Tracing where simplicity is actually used shows the caveat is
unnecessary.

Simplicity enters the reduction exactly once, in Lemma~\ref{lem:dual-cubic},
and only to conclude that the dual is bridgeless: a dual bridge corresponds
to a primal \emph{loop}, and simple graphs have none. What is needed is
therefore looplessness, not simplicity. Multiple edges are harmless, since a
multiple edge is not a loop and creates no dual bridge.

Re-reading the construction with that in hand: if a face walk repeats a
vertex, joining $u_F$ to each \emph{occurrence} gives $u_F$ several edges to
that vertex --- multiple edges, not a loop. Each new face is still a triangle
$(u_F, v_i, v_{i+1})$, and $v_i \ne v_{i+1}$ because consecutive equality
would be a loop in $G$, which a simple $G$ does not have. So stellar
subdivision of a loopless plane graph yields a loopless plane triangulation,
with no connectivity hypothesis at all.

Two identities corroborate the construction before any of it is built. The
subdivision adds one vertex per face and one edge per dart-corner, so
$V' = V + F$, $E' = E + 2E = 3E$, and each old face of length $k$ becomes $k$
triangles, giving $F' = 2E$. Then
\[
  V' - E' + F' = (V + F) - 3E + 2E = V - E + F = 2 ,
\]
so Euler is preserved exactly; and a map all of whose faces are triangles
satisfies $3F' = 2E'$, which here reads $6E = 6E$. Both hold with no slack,
which is what a correct triangulating construction must give and what an
incorrect one generally does not.

What this route costs, then, is the construction itself and the verification
that its faces are triangles --- and the next subsection carries both out on
darts, so that even the phrase ``$u_F$ sits inside $F$'' is discharged rather
than appealed to. What it does not cost is any argument about
non-adjacent pairs, hence neither separation nor forbidden subgraphs. That is
a different kind of obligation from Routes~A and~B, not merely a smaller
one --- though, given how often the size of this residual has been misjudged
here, the comparison is offered for checking rather than as a conclusion.
\end{remark}

\subsection{Stellar subdivision on darts}

The construction above was described by placing a vertex inside a face. That
phrase is the only geometry in it, and it can be removed: the whole
subdivision is a substitution on darts, and the triangles come out of a
three-line computation.

\begin{definition}[stellar subdivision]\label{def:stellar}
\Sprose
Let $M = (D,\sigma,\alpha)$ be a map with face permutation
$\varphi = \sigma\alpha$. For each face orbit
$F = (d_0, d_1, \dots, d_{k-1})$, so that $\varphi(d_i) = d_{i+1}$ with
indices mod $k$, adjoin $2k$ fresh darts $a_0,\dots,a_{k-1}$ and
$b_0,\dots,b_{k-1}$, and set
\[
  \alpha'(a_i) = b_i, \qquad
  \sigma'(\alpha d_i) = a_i, \qquad
  \sigma'(a_i) = d_{i+1}, \qquad
  \sigma'(b_i) = b_{i-1},
\]
with $\alpha' = \alpha$ and $\sigma' = \sigma$ on all darts not named. Write
$M' = (D', \sigma', \alpha')$.
\end{definition}

That $\sigma'$ is again a permutation is exactly the identity
$\sigma(\alpha d_i) = \varphi(d_i) = d_{i+1}$: the new dart $a_i$ is spliced
into the rotation between $\alpha d_i$ and $d_{i+1}$, where the old rotation
went directly. Distinct faces name distinct darts $\alpha d_i$, since the
face orbits partition $D$, so the clauses do not collide.

\begin{theorem}[the subdivision triangulates]\label{thm:stellar-darts}
\Sverified{Mettapedia.GraphTheory.FourColor.GoertzelV24StellarEuler.orbitCount_stellarRotation}
\Sverified{Mettapedia.GraphTheory.FourColor.GoertzelV24StellarEuler.orbitCount_stellarFace_eq_two_mul}
\Sverified{Mettapedia.GraphTheory.FourColor.GoertzelV24StellarEuler.orbitCount_stellarFlip}
\Sverified{Mettapedia.GraphTheory.FourColor.GoertzelV24StellarEuler.stellar_euler}
\Sverified{Mettapedia.GraphTheory.FourColor.GoertzelV24StellarStructure.stellar_old_embedding}
\Sverified{Mettapedia.GraphTheory.FourColor.GoertzelV24StellarStructure.stellar_loopless}
\Sverified{Mettapedia.GraphTheory.FourColor.GoertzelV24StellarStructure.stellar_face_three_vertices}
\Sverified{Mettapedia.GraphTheory.FourColor.GoertzelV24StellarRotationSystem.orbitFacesTwoSided}
\Sverified{Mettapedia.GraphTheory.FourColor.GoertzelV24StellarRotationSystem.orbitFaceDarts_card_three}
\Sverified{Mettapedia.GraphTheory.FourColor.GoertzelV24StellarDualStructure.stellar_primalConnected}
\Sverified{Mettapedia.GraphTheory.FourColor.GoertzelV24StellarDualStructure.stellar_euler_semantic}
The faces of $M'$ are exactly the $3$-cycles
\[
  (\, d_i, \; a_i, \; b_{i-1} \,),
\]
one for each face $F$ of $M$ and each index $i$. Consequently
$V' = V + F$, $E' = 3E$, $F' = 2E$, and $M'$ has the same Euler
characteristic as $M$.  The underlying graph of $M$ is a subgraph of that of
$M'$, and connectedness of $M$ implies connectedness of $M'$.  If $M$ is
loopless then so is $M'$, and each face of
$M'$ meets three distinct vertices; moreover the two darts of every edge of
$M'$ belong to distinct face orbits.
\end{theorem}

\begin{proof}
Compute $\varphi' = \sigma'\alpha'$ on the three darts in question:
\[
  \varphi'(d_i) = \sigma'(\alpha d_i) = a_i, \qquad
  \varphi'(a_i) = \sigma'(b_i) = b_{i-1}, \qquad
  \varphi'(b_{i-1}) = \sigma'(a_{i-1}) = d_i .
\]
So each such triple is a $\varphi'$-orbit of length three. Every dart of $D'$
occurs in exactly one triple --- each $d_i$ in its own, each $a_i$ in the
$i$-th, each $b_{i-1}$ in the $i$-th --- so these are all the faces.

For the counts: the darts $b_0,\dots,b_{k-1}$ form a single $\sigma'$-cycle,
one new vertex per face, and no old vertex is split, so $V' = V + F$; the new
edges are the pairs $\{a_i,b_i\}$, one per dart of $M$, so
$E' = E + 2E = 3E$; and the triples number $\sum_F k_F = 2E$, so $F' = 2E$.
Then $V' - E' + F' = (V+F) - 3E + 2E = V - E + F$.

No old edge is removed.  Each new vertex is the $b$-cycle belonging to an old
face and is joined by its new edges to vertices incident with that nonempty
face orbit.  Hence the old underlying graph is a subgraph of the new one, and
every new vertex attaches to it; connectedness is therefore preserved.

Each new edge joins the vertex carrying $\alpha d_i$ to the new vertex of
$F$, which is not an old vertex, so no new loop appears; old edges are
unchanged. Finally the face $(d_i, a_i, b_{i-1})$ meets the tail of $d_i$,
the head of $d_i$ (since $a_i$ sits at the vertex of $\alpha d_i$), and the
new vertex of $F$; these are three distinct vertices precisely because $d_i$
is not a loop.

For two-sidedness, an old edge has distinct old darts $d$ and $\alpha d$, and
the displayed partition puts exactly one old dart in each new face, so its two
darts lie in distinct faces.  A new edge $\{a_i,b_i\}$ meets the faces indexed
by $d_i$ and $d_{i+1}$.  These are distinct because a face orbit of length one
would give $\sigma(\alpha d_i)=d_i$, placing both darts of the old edge at one
vertex and making it a loop.  Thus every edge of $M'$ has two distinct sides.
\end{proof}

\begin{remark}[what is left of the embedding]
\Sprose
Definition~\ref{def:stellar} names four values of two permutations; the proof
evaluates $\sigma'\alpha'$ three times. No face interior, no curve, and no
plane occurs in either. Since the Euler characteristic is preserved, a
spherical $M$ yields a spherical $M'$, which is what
Lemma~\ref{lem:triangulate} was wanted for --- with containment of $M$ in
place of the spanning condition, which is all the reduction uses, a proper
colouring restricting to a subgraph regardless of whether the vertex set
grew.
\end{remark}

\begin{remark}[why the maximality lemma remains open]
\Sprose
Lemma~\ref{lem:triangulate} keeps its \textsc{open} marker because stellar
subdivision does not prove that stronger, same-vertex-set statement.  But it
is no longer a premise of Theorem~\ref{thm:tait-reduction}.  The reduction
needs a loopless triangulating supermap containing $G$, not a maximal simple
supergraph on exactly the old vertex set, and
Theorem~\ref{thm:stellar-darts} supplies precisely the weaker object.  This is
a change of lemma, not a proof of the old lemma, and the status labels preserve
that distinction.
\end{remark}

\begin{remark}[the shape of the debt]
\Sprose
The combinatorial-map reduction needs neither new algebra nor a Jordan curve theorem.  At the
ordinary-mathematics level it consists of explicit finite permutation
constructions plus the already-stated cycle-space theorem.  The stellar
construction, dual transport, literal-edge Klein-group potential, and
restriction to old vertices are now checked.  Translating a topological plane
drawing into this finite representation is a separate representation theorem,
recorded explicitly in Lemma~\ref{lem:planar-spherical-presentation}; it does
not weaken the kernel statement about finite combinatorial maps.
\end{remark}

\begin{remark}[how this closes the classical debts]
With Corollary~\ref{cor:two-sided} established, the route's internal
graph-backed presentation and Theorem~\ref{thm:tait-reduction} are checked,
including the least-counterexample wrapper of
Theorem~\ref{thm:spherical-headline-from-no-minimal}.  The standard finite
plane-embedding representation theorem remains named only in the separate
topological Corollary~\ref{cor:tait-reduction-topological}; keeping that
boundary explicit prevents the refuted legacy planarity predicate from
re-entering the kernel headline.
\end{remark}
